\chapter{出错处理设计}

\section{错误类型}
LingoBridge 平台在跨国互联网环境和多租户试点模式下运行时，可能会遭遇各层级的软件及网络硬件故障。为了确保系统在出现异常时不会发生全局崩溃或静默丢失学生数据，系统对常见错误进行了系统性的分类捕获与编码。系统典型的错误类型及协同处理矩阵如表 7.1 所示。

\begin{table}[H]
\centering
\caption{LingoBridge 系统核心错误分类与容错处理矩阵}
\label{tab:err_matrix}
\begin{tabularx}{\textwidth}{l p{2.2cm} p{3.2cm} X}
\toprule
\textbf{错误分类} & \textbf{网络状态码} & \textbf{错误特征描述} & \textbf{系统协同容错与优雅捕获机制} \\
\midrule
鉴权凭证失效 & 401 & 用户的 JWT 令牌过期或签名被篡改。 & 拦截器拦截请求，前端自动保留当前草稿并重定向至登录门户。 \\
课件体积超限 & 413 & 上传的 PDF/XLSX 课件文件体积超过 20MB。 & 前端在上传前静态计算字节数并拦截，避免产生无效跨海传输。 \\
持久层故障 & 500 & 关系型数据库死锁或连接池爆满。 & 仓库层自动回滚当前事务，向客户端返回友好错误，保障数据一致性。 \\
信道网络断连 & 1006 & 跨海网络抖动导致 WebSocket 连接中断。 & 前端自动进入无声心跳监测，后台启动指数退避重连算法。 \\
\bottomrule
\end{tabularx}
\end{table}

\section{异常处理策略}
针对上述列出的典型出错场景，平台在前后端设计了三层拦截与防御机制。

在最外层即浏览器用户界面层，所有的请求动作均被包装在防抖和节流拦截器中，当用户连续点击或网络响应滞后时，界面关键按钮自动置为禁用状态，杜绝高频请求堆积；同时，全局错误监听器会捕获所有未处理的 Promise 拒绝事件，并将其转化为用户友好的三语提示。

在中间控制层即 Express 网关路由处，所有的控制器均使用标准的 \texttt{try-catch} 代码块包裹，当发生类型越界、外键关联冲突等异常时，拦截器会统一格式化输出 JSON 错误包，包含标准的 \texttt{success: false} 属性以及精准的故障描述码。

在最深层即持久层，Repository 执行更新操作时，严格执行事务完整性控制（ACID），若检测到任何一步执行失败，立即在连接池中派发 \texttt{ROLLBACK} 命令，保证数据库状态的一致性。

\section{日志记录机制}
平台在后端内置了轻量级、高内聚的异步日志审计模块，精准记录高价值安全故障。

针对教师导入 Excel 作业和上传课件的核心通道，系统在 \texttt{assignment\_imports} 实体中专设了日志记录字段，每次导入操作不论成败，其解析过的任务计数、重复记录数量、错误字段所在的 Excel 行号以及处理时间戳都会被自动组装为紧凑的 JSON 数据结构持久化至物理表中。

在常规请求拦截维度，拦截器会自动记录包含用户 UUID、请求 IP 归属、目标 API 路径、HTTP 状态码以及故障源信息的单行结构化日志，为后期的安全审计和漏洞回溯提供核心物证。

\section{故障恢复与自愈逻辑}
为了确保在弱网环境下教学体验的连贯性，系统在前端信令层引入了基于指数退避的长连接弹性自愈架构。

当跨海网络抖动导致 WebSocket 信道异常断连（如捕获到 \texttt{close} 事件且状态码为 1006）时，客户端前端并不直接弹出阻断式错误，而是自动激活在后台静默运行的自适应指数退避与重连自愈逻辑。重连逻辑由基础延迟常量（设为 1000 毫秒）与最大重试次数限制（设为 10 次）驱动。每一次重连重试都会根据重试次数指数级递增等待时间，并加入随机抖动因子（Jitter，通常为 0 到 500 毫秒之间），以防多客户端因心跳闪断在同一瞬间向服务端发起高频握手握手，导致服务器端遭遇连接雪崩故障。

当重连重试次数未达上限且与后端 WSS 重新握手成功时，系统会在无声中恢复 Presence 在线状态同步并清空重试计数器。若重试次数突破最大上限，系统则判定连接彻底超时，优雅退化至离线安全降级模式，并在前台弹出非阻断式警告，重定向至登录门户以保障用户会话安全。底层的具体重连指数算法伪代码和物理执行序列下沉至详细设计说明书（LLD）中。
